\documentclass[11pt,a4paper]{report}

\usepackage[utf8]{inputenc}
\usepackage[T1]{fontenc}
\usepackage{lmodern}
\usepackage[margin=1in]{geometry}
\usepackage{graphicx}
\usepackage{xcolor}
\usepackage{listings}
\usepackage{hyperref}
\usepackage{booktabs}
\usepackage{enumitem}
\usepackage{fancyhdr}
\usepackage{titlesec}
\usepackage{tocloft}
\usepackage{skak}           % Chess diagrams and notation
\usepackage{xskak}          % Extended chess support
\usepackage{chessboard}     % Chess board diagrams
\usepackage{amsmath}
\usepackage{algorithm}
\usepackage{algpseudocode}

% Colors
\definecolor{codebackground}{RGB}{245,245,245}
\definecolor{codecomment}{RGB}{0,128,0}
\definecolor{codekeyword}{RGB}{0,0,180}
\definecolor{codestring}{RGB}{163,21,21}
\definecolor{linkcolor}{RGB}{0,80,160}
\definecolor{chaptercolor}{RGB}{40,60,100}

% Hyperref setup
\hypersetup{
    colorlinks=true,
    linkcolor=linkcolor,
    urlcolor=linkcolor,
    citecolor=linkcolor,
    pdftitle={Kirin Chess Engine Documentation},
    pdfauthor={Strydr Silverberg},
}

% Code listing style
\lstdefinestyle{codestyle}{
    backgroundcolor=\color{codebackground},
    basicstyle=\ttfamily\small,
    breakatwhitespace=false,
    breaklines=true,
    captionpos=b,
    commentstyle=\color{codecomment}\itshape,
    keywordstyle=\color{codekeyword}\bfseries,
    stringstyle=\color{codestring},
    numberstyle=\tiny\color{gray},
    numbers=left,
    numbersep=8pt,
    showspaces=false,
    showstringspaces=false,
    showtabs=false,
    tabsize=4,
    frame=single,
    framerule=0.5pt,
    rulecolor=\color{gray!50},
}
\lstset{style=codestyle}

% Chapter and section styling
\titleformat{\chapter}[display]
    {\normalfont\huge\bfseries\color{chaptercolor}}
    {\chaptertitlename\ \thechapter}{20pt}{\Huge}
\titlespacing*{\chapter}{0pt}{0pt}{30pt}

% Header and footer
\pagestyle{fancy}
\fancyhf{}
\fancyhead[L]{\leftmark}
\fancyhead[R]{\thepage}
\setlength{\headheight}{14pt}
\renewcommand{\headrulewidth}{0.4pt}

% Chess move formatting
\newcommand{\move}[1]{\texttt{#1}}

% Function/method documentation
\newcommand{\function}[4]{%
    \paragraph{\texttt{#1}} \hfill \textit{#2}\\
    \textbf{Parameters:} #3\\
    \textbf{Returns:} #4
}

% API endpoint
\newcommand{\apiendpoint}[3]{%
    \noindent\fbox{\texttt{#1}} \texttt{#2}\\
    #3\vspace{0.5em}
}

% Note box
\newcommand{\notebox}[1]{%
    \begin{center}
    \fcolorbox{gray}{gray!10}{%
        \parbox{0.9\textwidth}{\textbf{Note:} #1}
    }
    \end{center}
}

% Warning box
\newcommand{\warningbox}[1]{%
    \begin{center}
    \fcolorbox{red!70}{red!10}{%
        \parbox{0.9\textwidth}{\textbf{Warning:} #1}
    }
    \end{center}
}

%------------------------------------------------------------------------------
% Document Information
%------------------------------------------------------------------------------
\title{%
    \vspace{-1cm}
    \rule{\textwidth}{1pt}\\[0.5cm]
    {\Huge\bfseries Kirin}\\[0.3cm]
    {\Large Chess Engine Documentation}\\[0.3cm]
    \rule{\textwidth}{1pt}\\[1cm]
    {\large Current Repository Build}
}
\author{Strydr Silverberg\\strydr\_silverberg@mines.edu\\\\Colorado School of Mines}
\date{\today}

\begin{document}

\maketitle
\thispagestyle{empty}

\vfill
\begin{center}
    \textit{Kirin: A high-performance chess engine written in C++ powering an autonomous chess system}\\[1em]
    \textbf{2026 Colorado School of Mines ProtoFund Recipient}\\
    Team: Strydr Silverberg, Colin Dake
\end{center}
\vfill

\newpage
\tableofcontents
\newpage

%==============================================================================
\chapter{Introduction}
%==============================================================================

\section{Overview}

This document provides comprehensive documentation for the \textbf{Kirin Chess Engine}. Kirin is a chess engine written in C++ and an integral part of the \textbf{Kirin} autonomous chess system.

Kirin is the computational engine powering an autonomous chess system designed to deliver a complete single-player experience. By integrating intelligent move generation with a physical board capable of independent piece manipulation, the system eliminates the need for a human opponent while preserving the tactile satisfaction of over-the-board play.

\section{Features}

\begin{itemize}
    \item \textbf{Bitboard Representation}: Efficient 64-bit integer board representation for fast move generation and evaluation
    \item \textbf{Magic Bitboards}: Pre-computed attack tables for sliding pieces (bishops and rooks) using magic number multiplication
    \item \textbf{Zobrist Hashing}: Position hashing for transposition table lookups and repetition detection
    \item \textbf{Alpha-Beta Search}: Negamax search with alpha-beta pruning and iterative deepening
    \item \textbf{UCI Protocol}: Full Universal Chess Interface compatibility for GUI integration
    \item \textbf{Advanced Evaluation}: Material counting, piece-square tables, pawn structure analysis
    \item \textbf{Move Ordering}: MVV-LVA, killer moves, and history heuristic for search efficiency
    \item \textbf{Time Management}: Configurable time controls with lag compensation
\end{itemize}

\section{System Requirements}

\begin{table}[h]
\centering
\begin{tabular}{ll}
\toprule
\textbf{Component} & \textbf{Requirement} \\
\midrule
Operating System & Linux (primary development target) \\
Memory & Minimum 256 MB RAM \\
Processor & Any modern x86-64 CPU \\
Compiler & GCC or Clang with C++17 support \\
Dependencies & CMake 3.10+, POSIX headers (\texttt{unistd.h}, \texttt{sys/time.h}) \\
\bottomrule
\end{tabular}
\caption{Minimum system requirements}
\end{table}

%==============================================================================
\chapter{Installation}
%==============================================================================

\section{Building from Source}

Clone the repository and build:

\begin{lstlisting}[language=bash]
git clone git@github.com:strvdr/kirin-autonomous-chess.git
cd kirin-autonomous-chess/kirin
cmake -S . -B build
cmake --build build
./build/kirin
\end{lstlisting}

\section{Pre-built Binaries}

Build from source is the primary supported workflow for the current repository state.

\section{Start Playing}
Kirin is a UCI compatible chess engine, so configuring something like \href{https://github.com/lichess-bot-devs/lichess-bot}{Lichess Bot} is very simple. I have also created a chess CLI tool for viewing and playing games called \href{https://github.com/strvdr/zduel}{zduel} which you can use to play games against Kirin from the terminal.\\

If you want to play around with the engine but not an actual game, run the compiled binary and try this:
\begin{lstlisting}[language=bash]
kirin >> id name Kirin
kirin >> author name Strydr Silverberg
kirin >> uciok

user >> ucinewgame
user >> isready

kirin >> readyok

user >> go depth 12
\end{lstlisting}

%==============================================================================
\chapter{Architecture}
%==============================================================================

\section{Overview}

Kirin uses a bitboard-based architecture optimized for modern 64-bit processors. The engine is structured around several core components: board representation, move generation, position evaluation, and tree search.

\section{Board Representation}

\subsection{Bitboard Data Structures}

The board is represented using 12 bitboards (one per piece type and color) plus 3 occupancy bitboards:

\begin{lstlisting}[language=C++]
// Define Bitboard Type
#define U64 unsigned long long

// Piece bitboards: P, N, B, R, Q, K, p, n, b, r, q, k
U64 bitboards[12]; 

// Occupancy bitboards: white, black, both
U64 occupancy[3];

// Game state
int side;              // Side to move (white=0, black=1)
int enpassant;         // En passant target square
int castle;            // Castling rights (4 bits)
U64 hashKey;           // Zobrist hash of position
\end{lstlisting}

\subsection{Square Encoding}

Squares are encoded as integers 0-63, with a8=0 and h1=63:

\begin{lstlisting}[language=C++]
enum {
  a8, b8, c8, d8, e8, f8, g8, h8,
  a7, b7, c7, d7, e7, f7, g7, h7,
  a6, b6, c6, d6, e6, f6, g6, h6,
  a5, b5, c5, d5, e5, f5, g5, h5,
  a4, b4, c4, d4, e4, f4, g4, h4,
  a3, b3, c3, d3, e3, f3, g3, h3,
  a2, b2, c2, d2, e2, f2, g2, h2,
  a1, b1, c1, d1, e1, f1, g1, h1, no_sq
};
\end{lstlisting}

\subsection{Piece Encoding}

Pieces are encoded as integers 0-11:

\begin{lstlisting}[language=C++]
enum { P, N, B, R, Q, K, p, n, b, r, q, k };
\end{lstlisting}

\subsection{Castling Rights}

Castling rights are stored as a 4-bit value:

\begin{lstlisting}[language=C++]
// Binary encoding:
// 0001 (1) = white king-side
// 0010 (2) = white queen-side
// 0100 (4) = black king-side
// 1000 (8) = black queen-side
enum { wk = 1, wq = 2, bk = 4, bq = 8 };
\end{lstlisting}

\subsection{Bit Manipulation Macros}

\begin{lstlisting}[language=C++]
#define setBit(bitboard, square) ((bitboard) |= (1ULL << (square)))
#define getBit(bitboard, square) ((bitboard) & (1ULL << (square)))
#define popBit(bitboard, square) ((bitboard) &= ~(1ULL << (square)))
\end{lstlisting}

\subsection{Chess Position Example}

Here is an example starting position:

\begin{center}
\chessboard[setfen=rnbqkbnr/pppppppp/8/8/8/8/PPPPPPPP/RNBQKBNR w KQkq - 0 1]
\end{center}

After \move{1.e4 e5 2.Nf3}:

\begin{center}
\chessboard[setfen=rnbqkbnr/pppp1ppp/8/4p3/4P3/5N2/PPPP1PPP/RNBQKB1R b KQkq - 1 2]
\end{center}

\section{Move Generation}

\subsection{Move Encoding}

Moves are encoded as 24-bit integers containing all move information:

\begin{lstlisting}[language=C++]
// Binary move representation:
// 0000 0000 0000 0000 0011 1111  source square      0x3f
// 0000 0000 0000 1111 1100 0000  target square      0xfc0
// 0000 0000 1111 0000 0000 0000  piece              0xf000
// 0000 1111 0000 0000 0000 0000  promoted piece     0xf0000
// 0001 0000 0000 0000 0000 0000  capture flag       0x100000
// 0010 0000 0000 0000 0000 0000  double pawn push   0x200000
// 0100 0000 0000 0000 0000 0000  en passant flag    0x400000
// 1000 0000 0000 0000 0000 0000  castling flag      0x800000

#define encodeMove(source, target, piece, promoted, capture, dpp, enpassant, castle) \
  (source) | (target << 6) | (piece << 12) | (promoted << 16) | \
  (capture << 20) | (dpp << 21) | (enpassant << 22) | (castle << 23)
\end{lstlisting}

\subsection{Attack Tables}

Leaper pieces (pawns, knights, kings) use pre-computed attack tables:

\begin{lstlisting}[language=C++]
U64 pawn_attacks[2][64];   // [side][square]
U64 knightAttacks[64];
U64 kingAttacks[64];
\end{lstlisting}

\subsection{Magic Bitboards}

Sliding pieces (bishops, rooks, queens) use magic bitboard technique for O(1) attack generation:

\begin{lstlisting}[language=C++]
U64 bishopMasks[64];
U64 rookMasks[64];
U64 bishopAttacks[64][512];   // 32KB
U64 rookAttacks[64][4096];    // 256KB
U64 bishopMagicNums[64];
U64 rookMagicNums[64];

static inline U64 getBishopAttacks(int square, U64 occupancy) { 
  occupancy &= bishopMasks[square];
  occupancy *= bishopMagicNums[square];
  occupancy >>= 64 - bishopRelevantBits[square];
  return bishopAttacks[square][occupancy];
}
\end{lstlisting}

\subsection{Legal Move Generation}

\begin{algorithm}
\caption{Legal Move Generation}
\begin{algorithmic}[1]
\Procedure{GenerateMoves}{position}
    \State moves $\gets$ empty list
    \For{each piece type}
        \State bitboard $\gets$ pieces of this type
        \While{bitboard has pieces}
            \State square $\gets$ get least significant bit index
            \State Generate pseudo-legal moves for piece at square
            \For{each pseudo-legal move}
                \State Make move on board
                \If{own king not in check}
                    \State Add move to list
                \EndIf
                \State Restore board state
            \EndFor
            \State Remove bit from bitboard
        \EndWhile
    \EndFor
    \State \Return moves
\EndProcedure
\end{algorithmic}
\end{algorithm}

%==============================================================================
\chapter{Search Algorithm}
%==============================================================================

\section{Overview}

Kirin uses a negamax search with alpha-beta pruning, iterative deepening, and aspiration windows.

\section{Alpha-Beta Search}

\begin{algorithm}
\caption{Alpha-Beta with Negamax}
\begin{algorithmic}[1]
\Function{AlphaBeta}{position, depth, $\alpha$, $\beta$}
    \If{depth = 0}
        \State \Return \Call{QuiescenceSearch}{position, $\alpha$, $\beta$}
    \EndIf
    \If{position is repetition or fifty-move draw}
        \State \Return 0
    \EndIf
    \State Check transposition table for cached result
    \For{each move in \Call{GenerateMoves}{position}}
        \State \Call{MakeMove}{position, move}
        \State score $\gets$ $-$\Call{AlphaBeta}{position, depth$-$1, $-\beta$, $-\alpha$}
        \State \Call{UnmakeMove}{position, move}
        \If{score $\geq \beta$}
            \State Store killer move
            \State \Return $\beta$ \Comment{Beta cutoff}
        \EndIf
        \If{score $> \alpha$}
            \State $\alpha \gets$ score
            \State Update history heuristic
            \State Update PV table
        \EndIf
    \EndFor
    \State Store position in transposition table
    \State \Return $\alpha$
\EndFunction
\end{algorithmic}
\end{algorithm}

\section{Iterative Deepening}

The engine uses iterative deepening with aspiration windows:

\begin{lstlisting}[language=C++]
void searchPosition(int depth) {
  int score = 0;
  int alpha = -infinity;
  int beta = infinity;
  
  for(int currentDepth = 1; currentDepth <= depth; currentDepth++) {
    score = negamax(alpha, beta, currentDepth);
    
    // Aspiration window for next iteration
    alpha = score - 50;
    beta = score + 50;
    
    // Print UCI info
    printf("info score cp %d depth %d nodes %ld time %d pv ...", 
           score, currentDepth, nodes, getTimeMS() - startTime);
  }
  printf("bestmove ...");
}
\end{lstlisting}

\section{Search Enhancements}

\subsection{Transposition Table}

Zobrist hashing is used for position identification:

\begin{lstlisting}[language=C++]
U64 pieceKeys[12][64];    // Random keys for each piece on each square
U64 enpassantKeys[64];    // Random keys for en passant squares
U64 castlingKeys[16];     // Random keys for castling states
U64 sideKey;              // Random key for side to move

U64 generateHashKey() { 
  U64 finalKey = 0ULL; 
  for(int piece = P; piece <= k; piece++) { 
    U64 bitboard = bitboards[piece];
    while(bitboard) { 
      int square = getLSBindex(bitboard);
      finalKey ^= pieceKeys[piece][square];
      popBit(bitboard, square);
    }
  }
  if(enpassant != no_sq) finalKey ^= enpassantKeys[enpassant];
  finalKey ^= castlingKeys[castle];
  if(side == black) finalKey ^= sideKey;
  return finalKey;
}
\end{lstlisting}

\subsection{Move Ordering}

Moves are ordered for optimal alpha-beta pruning:

\begin{enumerate}
    \item Hash move (from transposition table)
    \item Winning captures (MVV-LVA ordering)
    \item Killer moves (quiet moves that caused beta cutoffs)
    \item History heuristic (successful quiet moves)
    \item Remaining quiet moves
\end{enumerate}

\subsubsection{MVV-LVA Table}

Most Valuable Victim - Least Valuable Attacker scoring:

\begin{lstlisting}[language=C++]
// MVV LVA [attacker][victim]
static int mvvLVA[12][12] = {
  105, 205, 305, 405, 505, 605, ...  // Pawn captures
  104, 204, 304, 404, 504, 604, ...  // Knight captures
  103, 203, 303, 403, 503, 603, ...  // Bishop captures
  102, 202, 302, 402, 502, 602, ...  // Rook captures
  101, 201, 301, 401, 501, 601, ...  // Queen captures
  100, 200, 300, 400, 500, 600, ...  // King captures
};
\end{lstlisting}

\subsection{Repetition Detection}

The engine maintains a repetition table to detect threefold repetition:

\begin{lstlisting}[language=C++]
U64 repetitionTable[1000]; 
int repetitionIndex;

static inline int isRepetition() {
  for(int i = 0; i < repetitionIndex; i++) {
    if(repetitionTable[i] == hashKey) return 1;
  }
  return 0;
}
\end{lstlisting}

%==============================================================================
\chapter{Evaluation}
%==============================================================================

\section{Overview}

The evaluation function returns a score in centipawns from the perspective of the side to move. Positive scores favor the side to move.

\section{Material}

\begin{table}[h]
\centering
\begin{tabular}{lc}
\toprule
\textbf{Piece} & \textbf{Value (centipawns)} \\
\midrule
Pawn & 100 \\
Knight & 300 \\
Bishop & 350 \\
Rook & 500 \\
Queen & 1000 \\
King & 10000 \\
\bottomrule
\end{tabular}
\caption{Piece values in Kirin}
\end{table}

\section{Positional Factors}

\subsection{Piece-Square Tables}

Each piece type has a positional bonus table. Example for pawns:

\begin{lstlisting}[language=C++]
const int pawnScore[64] = {
    90,  90,  90,  90,  90,  90,  90,  90,
    30,  30,  30,  40,  40,  30,  30,  30,
    20,  20,  20,  30,  30,  30,  20,  20,
    10,  10,  10,  20,  20,  10,  10,  10,
     5,   5,  10,  20,  20,   5,   5,   5,
     0,   0,   0,   5,   5,   0,   0,   0,
     0,   0,   0, -10, -10,   0,   0,   0,
     0,   0,   0,   0,   0,   0,   0,   0
};
\end{lstlisting}

\subsection{Knight Positioning}

Knights are rewarded for central placement:

\begin{lstlisting}[language=C++]
const int knightScore[64] = {
    -5,   0,   0,   0,   0,   0,   0,  -5,
    -5,   0,   0,  10,  10,   0,   0,  -5,
    -5,   5,  20,  20,  20,  20,   5,  -5,
    -5,  10,  20,  30,  30,  20,  10,  -5,
    -5,  10,  20,  30,  30,  20,  10,  -5,
    -5,   5,  20,  10,  10,  20,   5,  -5,
    -5,   0,   0,   0,   0,   0,   0,  -5,
    -5, -10,   0,   0,   0,   0, -10,  -5
};
\end{lstlisting}

\section{Pawn Structure}

\subsection{Double Pawns}

Doubled pawns receive a penalty:

\begin{lstlisting}[language=C++]
const int doublePawnPenalty = -10;

doublePawns = countBits(bitboards[P] & fileMasks[square]);
if(doublePawns > 1) score += doublePawns * doublePawnPenalty;
\end{lstlisting}

\subsection{Isolated Pawns}

Pawns with no friendly pawns on adjacent files are penalized:

\begin{lstlisting}[language=C++]
const int isolatedPawnPenalty = -10;

if((bitboards[P] & isolatedMasks[square]) == 0) 
    score += isolatedPawnPenalty;
\end{lstlisting}

\subsection{Passed Pawns}

Passed pawns receive bonuses based on how advanced they are:

\begin{lstlisting}[language=C++]
const int passedPawnBonus[8] = { 0, 10, 30, 50, 75, 100, 150, 200 };

if((passedWhiteMasks[square] & bitboards[p]) == 0) 
    score += passedPawnBonus[getRank[square]];
\end{lstlisting}

%==============================================================================
\chapter{UCI Protocol}
%==============================================================================

\section{Overview}

The engine communicates via the Universal Chess Interface (UCI) protocol.

\section{Supported Commands}

\subsection{Standard UCI Commands}

\apiendpoint{IN}{uci}{Initialize UCI mode and request engine identification.}

\apiendpoint{IN}{isready}{Check if engine is ready. Engine responds with \texttt{readyok}.}

\apiendpoint{IN}{ucinewgame}{Reset engine state for a new game.}

\apiendpoint{IN}{position [fen | startpos] moves ...}{Set up a position.}

\apiendpoint{IN}{go [options]}{Start calculating. Options include:}
\begin{itemize}[noitemsep]
    \item \texttt{depth n} -- search to depth n
    \item \texttt{movetime ms} -- search for ms milliseconds
    \item \texttt{wtime/btime ms} -- time remaining for white/black
    \item \texttt{winc/binc ms} -- increment per move
    \item \texttt{movestogo n} -- moves until next time control
    \item \texttt{infinite} -- search until stopped
    \item \texttt{perft n} -- run perft test to depth n
\end{itemize}

\apiendpoint{IN}{stop}{Stop calculating.}

\apiendpoint{IN}{quit}{Exit the engine.}

\subsection{Engine Responses}

\begin{lstlisting}
id name Kirin
id author Strydr Silverberg
uciok

info score cp 35 depth 8 nodes 234567 time 1250 pv e2e4 e7e5 g1f3

bestmove e2e4
\end{lstlisting}

%==============================================================================
\section{Skill Level}
%==============================================================================

Kirin exposes a three-tier difficulty setting that adjusts both search depth and
move quality, making the engine accessible to players of all levels while
preserving full strength for competitive play.

\subsection{Difficulty Tiers}

\begin{table}[h]
\centering
\begin{tabular}{clll}
\toprule
\textbf{Level} & \textbf{Name}   & \textbf{Max Search Depth} & \textbf{Move Noise} \\
\midrule
0              & Easy            & 3 plies                   & $\pm$150 centipawns \\
1              & Medium          & 5 plies                   & $\pm$50 centipawns  \\
2              & Hard (default)  & Caller-supplied depth      & None                \\
\bottomrule
\end{tabular}
\caption{Skill level configuration}
\end{table}

Two independent mechanisms are combined to simulate human-like play at lower
difficulty settings:

\begin{enumerate}
    \item \textbf{Depth capping.} Inside \texttt{searchPosition}, the
    caller-supplied depth is silently clamped to a tier-specific ceiling before
    iterative deepening begins. At skill~0 the engine never searches deeper than
    3 plies; at skill~1 the cap is 5 plies; at skill~2 the cap is set to 64
    (effectively unlimited for practical time controls).

    \item \textbf{Root-move noise.} After the search completes, the engine's
    best move is not always played. At skill levels below 2, each legal root
    move is assigned the engine's score perturbed by a random offset drawn
    uniformly from $[-\textit{range},\, +\textit{range}]$, where
    \textit{range} is 150\,cp at skill~0 and 50\,cp at skill~1. The move with
    the highest \emph{noisy} score is selected. This causes the engine to
    occasionally prefer a suboptimal move in a way that resembles human error
    rather than purely random play.
\end{enumerate}

\subsection{UCI Configuration}

The skill level is exposed as a standard UCI spin option and can be set from
any compatible GUI or from the command line:

\begin{lstlisting}[language=bash]
setoption name Skill Level value 0   % Easy
setoption name Skill Level value 1   % Medium
setoption name Skill Level value 2   % Hard (default)
\end{lstlisting}

The engine announces the option during UCI initialisation:

\begin{lstlisting}
option name Skill Level type spin default 2 min 0 max 2
\end{lstlisting}

Values outside the range [0,\,2] are clamped silently. The engine acknowledges
each change with an info string:

\begin{lstlisting}
info string Skill Level set to 1
\end{lstlisting}

\subsection{Implementation}

The global variable \texttt{skillLevel} in \texttt{search.cpp} stores the
current setting. It is read in two places:

\begin{lstlisting}[language=C++]
// Depth cap applied at the start of searchPosition()
static const int skillDepthCap[3] = { 3, 5, 64 };
int depthCap = skillDepthCap[skillLevel];
if (depth > depthCap) depth = depthCap;

// Noise injection applied after the best move is found
if (skillLevel < 2) {
    static const int noiseRange[2] = { 150, 50 };
    int range = noiseRange[skillLevel];
    // ... re-score root moves with random perturbation
}
\end{lstlisting}

\notebox{Skill level does not affect the correctness of move generation or
position evaluation — illegal moves are never produced regardless of the
difficulty setting. The engine test suite (\texttt{engine\_test.cpp}) verifies
this invariant across all three levels and multiple board positions.}
%==============================================================================
\chapter{API Reference}
%==============================================================================

\section{Core Functions}

\subsection{Board Management}

\function{parseFEN(const char *fen)}{Function}{fen -- FEN string describing position}{void}

Parses a FEN string and sets up the board position, including piece placement, side to move, castling rights, en passant square, and initializes the hash key.

\function{printBoard()}{Function}{None}{void}

Prints the current board position to stdout with piece symbols, side to move, castling rights, en passant square, and hash key.

\subsection{Move Generation}

\function{generateMoves(moves *moveList)}{Function}{moveList -- pointer to move list structure}{void}

Generates all pseudo-legal moves for the current position and stores them in the provided move list.

\function{makeMove(int move, int moveFlag)}{Function}{move -- encoded move, moveFlag -- allMoves or onlyCaptures}{int -- 1 if legal, 0 if illegal}

Attempts to make a move on the board. Returns 1 if the move is legal (doesn't leave own king in check), 0 otherwise.

\subsection{Search}

\function{searchPosition(int depth)}{Function}{depth -- maximum search depth}{void}

Searches the current position using iterative deepening and prints the best move via UCI protocol.

\function{negamax(int alpha, int beta, int depth)}{Function}{alpha, beta -- search window, depth -- remaining depth}{int -- position score}

Core search function implementing negamax with alpha-beta pruning.

\subsection{Evaluation}

\function{evaluate()}{Function}{None}{int -- position score in centipawns}

Returns a static evaluation of the current position from the perspective of the side to move.

\function{isAttacked(int square, int side)}{Function}{square -- target square, side -- attacking side}{int -- 1 if attacked, 0 otherwise}

Checks if a square is attacked by the specified side.

%==============================================================================
\chapter{Testing}
%==============================================================================

\section{Perft Testing}

Perft (performance test) validates move generation by counting all leaf nodes at a given depth:

\begin{lstlisting}[language=bash]
position startpos
go perft 6
\end{lstlisting}

\begin{table}[h]
\centering
\begin{tabular}{cr}
\toprule
\textbf{Depth} & \textbf{Nodes} \\
\midrule
1 & 20 \\
2 & 400 \\
3 & 8,902 \\
4 & 197,281 \\
5 & 4,865,609 \\
6 & 119,060,324 \\
\bottomrule
\end{tabular}
\caption{Perft results from starting position}
\end{table}

\section{Test Positions}

The engine includes several built-in test positions:

\begin{lstlisting}[language=C++]
#define startPosition "rnbqkbnr/pppppppp/8/8/8/8/PPPPPPPP/RNBQKBNR w KQkq - 0 1"
#define tricky_position "r3k2r/p1ppqpb1/bn2pnp1/3PN3/1p2P3/2N2Q1p/PPPBBPPP/R3K2R w KQkq - 0 1"
#define killer_position "rnbqkb1r/pp1p1pPp/8/2p1pP2/1P1P4/3P3P/P1P1P3/RNBQKBNR w KQkq e6 0 1"
#define cmk_position "r2q1rk1/ppp2ppp/2n1bn2/2b1p3/3pP3/3P1NPP/PPP1NPB1/R1BQ1RK1 b - - 0 9"
\end{lstlisting}

%==============================================================================
\chapter{Performance}
%==============================================================================

\section{Optimization Techniques}

\begin{itemize}
    \item \textbf{Magic bitboards}: O(1) sliding piece attack generation using pre-computed magic numbers
    \item \textbf{Incremental hash updates}: Zobrist hash is updated incrementally during make/unmake move
    \item \textbf{Copy-make}: Board state is copied before making moves for efficient unmake
    \item \textbf{Inline functions}: Critical functions are marked \texttt{static inline} for compiler optimization
    \item \textbf{Bit manipulation intrinsics}: Uses efficient bit counting and LSB detection
\end{itemize}

\section{Memory Usage}

\begin{table}[h]
\centering
\begin{tabular}{lr}
\toprule
\textbf{Component} & \textbf{Size} \\
\midrule
Bishop attack table & 32 KB \\
Rook attack table & 256 KB \\
Transposition table & Configurable \\
Repetition table & 8 KB \\
\bottomrule
\end{tabular}
\caption{Memory usage of major components}
\end{table}

%==============================================================================
\chapter{Contributing}
%==============================================================================

\section{Development Setup}

\begin{enumerate}
    \item Fork the repository
    \item Clone your fork locally
    \item Build with \texttt{make debug} for development builds
    \item Run tests with \texttt{make test}
\end{enumerate}

\section{Code Style}

\begin{itemize}
    \item Use camelCase for function and variable names
    \item Use SCREAMING\_CASE for macros and constants
    \item Keep functions focused and under 50 lines where possible
    \item Comment non-obvious algorithms and magic numbers
\end{itemize}

\section{Pull Request Process}

\begin{enumerate}
    \item Fork the repository
    \item Create a feature branch
    \item Make changes and verify perft results match expected values
    \item Submit a pull request with clear description
\end{enumerate}

%==============================================================================
\appendix
\chapter{Glossary}
%==============================================================================

\begin{description}
    \item[Alpha-Beta] A search algorithm that prunes branches that cannot affect the final result.
    \item[Bitboard] A 64-bit integer representing piece positions on the board, where each bit corresponds to a square.
    \item[Centipawn] One hundredth of a pawn; standard unit for evaluation scores.
    \item[FEN] Forsyth-Edwards Notation; a standard for describing chess positions in a single line of text.
    \item[Magic Bitboard] A technique using multiplication and lookup tables for O(1) sliding piece attack generation.
    \item[MVV-LVA] Most Valuable Victim - Least Valuable Attacker; a move ordering heuristic for captures.
    \item[Negamax] A variant of minimax that exploits the zero-sum property of chess.
    \item[Perft] Performance test; counts all leaf nodes at a given depth to validate move generation.
    \item[Ply] A half-move; one move by either white or black.
    \item[Quiescence Search] Extended search of capture sequences to avoid horizon effects.
    \item[UCI] Universal Chess Interface; standard protocol for chess engine communication.
    \item[Zobrist Hashing] A technique for hashing chess positions using XOR of random numbers.
    \item[Skill Level] A configurable difficulty setting (0--2) that limits search depth and injects random noise into move selection to simulate human-like play.
    \item[Move Noise] A random centipawn perturbation applied to root-move scores at lower skill levels, causing the engine to occasionally prefer a suboptimal move.
\end{description}

%==============================================================================
\chapter{References}
%==============================================================================

\begin{itemize}
    \item Chess Programming Wiki: \url{https://www.chessprogramming.org/}
    \item UCI Protocol Specification: \url{https://gist.github.com/DOBRO/2592c6dad754ba67e6dcaec8c90165bf}
    \item Stockfish source code: \url{https://github.com/official-stockfish/Stockfish}
    \item VICE Chess Engine by Richard Allbert (time control implementation reference)
    \item Tord Romstad's magic number generation proposal
    \item Maksim Korzh's chess programming guidance
\end{itemize}

%%%%%%%%%%%%%%%%%%%%%%%%%%%%%%%%%%%%%%%%%%%%%%%%%%%%%%%%%%%%%%%%%%%%%%%%%%%%%%%
\end{document}
%%%%%%%%%%%%%%%%%%%%%%%%%%%%%%%%%%%%%%%%%%%%%%%%%%%%%%%%%%%%%%%%%%%%%%%%%%%%%%%
